% 對應英文：sections/09_conclusion.tex

本文指出，加密貨幣橫截面報酬預測中，
深度學習表現不穩定的核心問題，
並不必然來自架構本身，而往往來自損失函數與評估目標的不一致。
當任務真正關心的是資產之間的相對排序時，
以原始報酬直接驅動 ListNet 訓練，容易因跨週波動尺度不同而產生不穩定梯度；
改用排名百分位目標，並結合梯度累積後，
深度學習模型的排序能力才開始穩定反映在樣本外績效上。

\paragraph{架構與訊號型態的對應關係。}
本文最重要的實證結論，不是單一模型全面勝出，
而是不同架構對不同訊號型態各有優勢。
CS-Gated 在鏈上特徵主導時表現最佳，
顯示當期橫截面基本面已足以支持高品質排序；
LSTM 與 TFT 則在納入川普社群媒體與 ETF 結構流量後明顯改善，
說明時序模型較能吸收事件性與延續性較強的替代資料。
同時，隨機森林在高維資訊集下仍具強勁競爭力，
提醒我們傳統非參數方法並未因深度學習的引入而失去參考價值。

\paragraph{方法論上的意涵。}
消融研究顯示，真正帶來躍升的不是單一新架構，
而是先重新設計監督訊號，再讓模型在較穩定的梯度條件下學習。
這個結果意味著，
過去文獻中所謂「深度學習不適合金融預測」的結論，
至少有一部分可能來自目標設定與投資任務不一致，而非模型容量不足。
對橫截面資產定價研究而言，
損失函數如何對應實際投資決策，往往與模型本身同等重要。

\paragraph{研究限制。}
本文仍有三項主要限制。
第一，測試期間僅 49 週，統計推論的精確度有限。
第二，資產宇宙以 2026 年 3 月市值固定選定，存在存活偏差。
第三，深度學習模型的最佳表現高度依賴多種子集成，
因此若從單次訓練穩定性或計算成本角度評估，
其優勢未必如最終集成績效所示般明顯。

\paragraph{未來方向。}
未來研究可沿三個方向延伸。
其一，將排名正規化方法應用到股票或其他橫截面資產市場，
檢驗這一方法論是否具有更廣泛的外部效度。
其二，納入更高頻的訂單簿、鏈上微結構或敘事型文本資料，
以區分哪些訊號真正需要時序編碼，哪些只需橫截面比較即可。
其三，除了十分位多空排序外，
亦可探索風險平價、約束式最佳化或動態倉位控制等投資組合建構方式，
使模型分數與最終交易決策之間的連結更加完整。
